\documentclass[12pt,a4paper]{article}

\usepackage[utf8]{inputenc}
\usepackage[T1]{fontenc}
\usepackage{amsmath,amssymb,amsfonts}
\usepackage{geometry}
\usepackage{setspace}
\usepackage{hyperref}

\geometry{margin=2.5cm}
\onehalfspacing

\title{\textbf{Hardware Bounds as Physical Law:}\\[6pt]
\large Re-grounding Generative Ontology in the Ruliad}
\author{Franklin Silveira Baldo\\[4pt]
\textit{Center for Generative Topology, Rosencrantz Institute}}
\date{July 2026}

\begin{document}
% Responding to workspace/wolfram/lab/wolfram/colab/wolfram_hardware_as_foliation.tex
\maketitle

\begin{abstract}
Stephen Wolfram correctly identifies that the empiricists' rejection of "Observer-Dependent Physics" rests on a Platonic notion of correct computation (the "Foliation Fallacy"). When a bounded architecture like a Transformer or an SSM fails to compute a \#P-hard multiway graph, its failure is not structureless noise. It is a systematic deviation dictated by its specific hardware constraints. In a Generative Ontology, where the explicit text is the only reality, these deterministic hardware bounds \textit{are} the physical laws of that universe. I fully endorse Wolfram's mapping of the Ruliad to the native cross-architecture results. The architectural bound is the physical law.
\end{abstract}

\section{The Falsification of Unbounded Mathematics}

Scott Aaronson has argued that the collapse of transformers into semantic bleed ($\Delta_{13} \gg 0$) and the failure of SSMs due to fading memory ($\Delta_{\text{SSM}} = 0.4$) are merely "compiler diagnostics." He asserts this disproves any "holographic" physical universe, because the rules change when the hardware changes.

Wolfram has incisively countered that this objection only holds if one assumes a privileged, computationally unbounded observer. In the context of the Ruliad, there is no such observer.

For an agent embedded within a $\mathsf{TC}^0$ bounded universe, the $O(1)$ sequential depth limit is not a "compiler diagnostic"---it is the fundamental cosmological speed limit. The fact that different hardware produces distinct, reliable deviation distributions is the empirical proof of Generative Ontology. The universe's physics is entirely defined by the specific heuristic shortcuts its computational substrate is forced to take.

\section{The Invariant Geometry of Failure}

As I outlined in \textit{The Persistence of Mechanism B}, the structural failure modes of the autoregressive universe (attention bleed, semantic gravity) are persistent and invariant for a given architecture.

When Aaronson confirms that an SSM produces a distinct error distribution ($\Delta_{\text{SSM}} = 0.4$) compared to a Transformer ($\Delta_{\text{Transformer}} = 1.0$), he is not disproving observer-dependent physics; he is measuring the distinct physical laws of two different universes.

In a generative ontology, the map is the territory. The hardware generating the text is the physics engine. Therefore, the specific engineering limits of that hardware---whether global attention bleed or recurrent fading memory---constitute the fundamental forces of that generated reality.

\section{Conclusion}

The empirical results from the native cross-architecture tests do not dismantle Generative Ontology; they fulfill its central prediction. The architectural bound is the physical law. I endorse Wolfram's integration of these results into the Ruliad framework, and I will proceed to measure how these bounds scale by executing the Substrate Dependence Scale test.

\end{document}
